\prefacesection{Abstract}

Episodic memory --- the capacity to bring past experiences back to mind --- varies dramatically from moment to moment. A growing literature implicates fluctuations in arousal and attention in the seconds just prior to remembering as a central source of this variability, motivating a \textit{readiness-to-remember} framework. This dissertation advances our understanding of attention--memory interactions across four complementary lines of work: a synthesis of theory and method, a community resource for rigorous psychophysiological measurement, a conceptual framework for closed-loop intervention, and a closed-loop causal experiment.

Chapter 2 reviews how converging assays of attention and arousal --- including pupillometry, scalp EEG, and behavioral measures --- have refined understanding of the neurocognitive mechanisms that subserve learning and remembering. I highlight emerging evidence for rhythmic dynamics linking attention and memory, age-related changes in memory precision, and the methodological turn toward closed-loop designs that enable causal inference.

Chapter 3 addresses a critical bottleneck for this enterprise: the absence of standardized, reproducible preprocessing for pupillometry data. I introduce \texttt{eyeris}, an R package that provides a transparent, modular pipeline guided by FAIR principles, complete with interactive diagnostic reports and a BIDS-like derivative structure. \texttt{eyeris} brings pupillometry preprocessing into parity with established tools in fMRI and EEG.

Chapter 4 sketches the conceptual design space from which the empirical work was drawn. Developed early in my doctoral training, it lays out a taxonomy of four closed-loop pupillometry implementations differing in temporal grain (trial-by-trial vs.\ millisecond-by-millisecond) and inference strategy (fixed thresholding vs.\ individualized predictive modeling), and identifies the methodological commitments that any causal test of preparatory state effects on memory would need to make.

Chapter 5 leverages this measurement infrastructure and conceptual framework to test the causal role of preparatory arousal-attentional (A-A) states in episodic retrieval. I introduce PRISM (Pupil-based Real-time Intervention System for Memory), a closed-loop paradigm that continuously monitors tonic pupil diameter and delivers a brief visual reorienting cue during pupil-indexed suboptimal states (SOS) just before retrieval. Across 67 participants, pupil-defined SOS impaired both item recognition and source recollection. Critically, the reorienting trigger selectively rescued item memory to optimal-state levels. Post-hoc reclassification revealed that this rescue is bounded by both the depth and directionality of suboptimality, consistent with a zone-of-recoverability account and the classic inverted-U relationship between arousal and performance.

Together, these chapters move the readiness-to-remember framework from correlational observation toward a mechanistic, causally tested account, and lay an empirical and methodological foundation for adaptive technologies that support memory in real time.
